% ALL NUMBERS from RESULTS_027_CAMERA_READY.md (camera-ready: 500 maps,
% density 0.27, stage-2 base, randomized attack, paired-bootstrap 95% CIs).
% Probe AUC now 500-map (0.99 oracle / 0.85-0.88 realistic, 0.88 at K_min=20)
% -- rerun done 2026-07-20, rerun/probe_{oracle,realistic}_500.txt; confirms
% the earlier 150-map feasibility figures. No dev-scale numbers remain.
% Verify each figure exists in manuscript/figures/ before final compile.

\subsection{Communication is load-bearing under sensor dropout}\label{sec:res-anchor}
We first establish that collaborative perception is not a refinement but a
requirement in this regime. Table~\ref{tab:anchor} is an
\emph{information ablation}: a single policy is evaluated on the same maps
with shared observations fused into its range channel (ON) and with its own
LiDAR only (OFF), under $\sim\!33\%$ sustained blindness at density $0.27$.
Because the weights are identical across the two arms, the gap is
attributable to the shared information alone. Sharing raises honest-drone
success from $45.9\%$ to $89.3\%$ and all-reach map success from $10.4\%$
to $67.8\%$ --- a $+57.4$ percentage-point map-level gap
(CI $[+52.8,+61.8]$). A policy \emph{trained} with sharing reaches a
similar $87.7\%$ in its native mode; we do not interpret the small
difference between the two, as it conflates two independently trained
policies. The consistent picture is that the benefit resides in the shared
information itself, which even a policy never trained with it exploits
zero-shot. A dropout sweep (Table~\ref{tab:dropout}) shows this
gap is an operating-point phenomenon: at $0\%$ dropout sharing confers no
advantage --- the gap is in fact slightly negative ($-1.3$ pp,
CI $[-2.3,-0.4]$), consistent with the early observation that cooperative
perception did not meaningfully improve over local
perception~\cite{trupercept2020} --- at $\sim\!33\%$ the gap opens to
$+41.4$ pp, and at $\sim\!50\%$ it widens further to $+50.8$ pp. We evaluate the attack and its defenses at the
$\sim\!33\%$ operating point, where collaboration is decisively valuable yet
the map remains solvable.

\begin{table}[t]\centering
\caption{Collaborative-perception anchor at density $0.27$, 500 maps.
Sharing (ON) versus own-LiDAR only (OFF) under $\sim\!33\%$ dropout.}
\label{tab:anchor}
\begin{tabular}{lccc}
\toprule
 & Drone-level & Map-level & Gap (map, 95\% CI) \\
\midrule
ON (shared)   & $89.3\%$ & $67.8\%$ & \multirow{2}{*}{$+57.4$ $[+52.8,+61.8]$} \\
OFF (own only)& $45.9\%$ & $10.4\%$ & \\
\bottomrule
\end{tabular}
\end{table}

\begin{table}[t]\centering
\caption{Dropout sweep, drone-level success, 500 maps. The comm advantage
appears only once sensing degrades.}
\label{tab:dropout}
\begin{tabular}{lcccc}
\toprule
Dropout & Blind (est.) & ON & OFF & Gap (95\% CI) \\
\midrule
$0\%$  & ${\sim}0\%$  & $88.7$ & $90.0$ & $-1.3$ $[-2.3,-0.4]$ \\
$10\%$ & ${\sim}33\%$ & $87.7$ & $46.3$ & $+41.4$ $[+38.9,+43.8]$ \\
$20\%$ & ${\sim}50\%$ & $86.2$ & $35.4$ & $+50.8$ $[+48.3,+53.3]$ \\
\bottomrule
\end{tabular}
\end{table}

\paragraph{Reading the baselines.}
Absolute success levels throughout this section should be read against the
deliberately stacked stresses. Every condition carries $\sim$33\% sustained
sensing blackout (the premise of collaboration), a calibrated obstacle
density chosen so the task is demanding but solvable, and --- as noise grows
--- position uncertainty approaching obstacle scale. The attack-free
ceiling therefore runs from $\sim$86\% at $\sigma=0$ down to $53.4\%$ at
$\sigma=0.6$; the latter is a navigation-information limit, not a defense
artifact (Section~\ref{sec:res-limit}). Crucially, the \emph{same} ceiling
underlies every arm of every comparison, so recovery and no-harm --- which
are paired differences against that common base --- remain fair measures of
the defense regardless of the absolute level.

\subsection{The minimum-fusion Byzantine vulnerability}\label{sec:res-vuln}
Because a fabricated near obstacle wins the per-ray minimum, even a small
number of traitors is damaging, and the camouflage placement is markedly
worse than the wall. Under attack with no defense, two traitors reduce
honest success from the $\sim\!86\%$ noiseless ceiling to $71.8\%$ (wall)
and $63.5\%$ (camouflage); three traitors reach $68.8\%$ and $60.9\%$
respectively. The damage \emph{saturates} between two and three traitors
(an added $\sim\!3$ pp), consistent with overlapping phantoms occupying the
same corridor. As anticipated by the fusion mechanism, the attack is
essentially independent of the dropout rate --- it overrides healthy and
degraded sensors alike --- so we hold dropout fixed and vary sensor noise
in what follows.

\subsection{Naive consistency trust is destructive under noise}\label{sec:res-naive}
Table~\ref{tab:naive} reports the naive single-frame filter against a
two-traitor wall attack as sensor noise increases. At $\sigma=0$ it is
excellent: precision $1.00$, and it recovers $+13.7$ pp. As soon as noise
appears it inverts. Detection precision collapses to $0.23$ --- more than
three quarters of its accusations are of honest neighbours --- and the
no-harm column goes deeply negative ($-27.5$ pp at $\sigma=0.4$), meaning
that \emph{with no attacker present at all}, merely enabling the filter
destroys a quarter of the swarm's success by gating out honest teammates.
Recovery under attack turns negative: the ``defense'' leaves the swarm worse
off than no defense. This is the direct empirical signature of the
$\sqrt{2}\sigma$ honest-disagreement problem of
Section~\ref{sec:methods-defense}.

\begin{table}[t]\centering
\caption{Naive single-frame trust, wall attack, $f=2$, 500 maps.
FP-harm $=$ no-harm cost; negative recovery means worse than no defense.}
\label{tab:naive}
\begin{tabular}{lccccc}
\toprule
$\sigma$ & Base & No-harm & Attack & Defense & Recovery (P/R)\\
\midrule
$0.0$ & $85.9$ & $85.9$ & $72.3$ & $85.9$ & $+13.7$ ($1.00/0.98$)\\
$0.2$ & $79.0$ & $65.1$ & $63.1$ & $62.8$ & $-0.3$ ($0.28/0.98$)\\
$0.4$ & $65.6$ & $38.1$ & $51.6$ & $38.3$ & $-13.3$ ($0.23/0.96$)\\
$0.6$ & $54.9$ & $33.3$ & $41.2$ & $33.1$ & $-8.1$ ($0.23/0.94$)\\
\bottomrule
\end{tabular}
\end{table}

\subsection{The robust filter restores safety but hits a camouflage limit}\label{sec:res-robust}
Widening the tolerance to $\varepsilon_0+4\sigma$ repairs the false
accusations: across the noise range the robust filter holds precision
$0.92$--$1.00$ and no-harm within $\pm0.3$ pp of base (the no-harm
confidence intervals span zero at every level). Against the conspicuous
wall attack it recovers gracefully. Against camouflage, however, it retains
a blind spot exactly where the paper's threat is sharpest: at $\sigma=0.6$,
$f=2$, its detection recall is only $0.13$ and its recovery is $+3.4$ pp
(CI $[+0.9,+5.8]$). The lie hidden inside the widened tolerance band cannot
be contradicted on any single frame. This is the limitation the temporal
test removes.

\subsection{Temporal offset-bias trust: the main result}\label{sec:res-temporal}
Table~\ref{tab:temporal} gives the full two-traitor results for both attack
placements across noise, comparing the single-frame robust filter with the
composed temporal filter. Three properties hold throughout. First, the
temporal filter never regresses: on the wall attack it matches or exceeds
the robust filter at every noise level. Second, its advantage opens exactly
where the robust filter fails --- at high noise under camouflage. At
$\sigma=0.6$, $f=2$, camouflage, temporal detection recall rises from
$0.13$ to $0.69$ and recovery from $+3.4$ to $+12.2$ pp
(CI $[+9.8,+14.9]$), restoring success from the attacked $37.6\%$ to
$49.8\%$ against a base of $53.4\%$. Third, this comes at no measurable cost
to honest swarms: the no-harm figure is $-0.4$ pp with a confidence
interval spanning zero.

\begin{table}[t]\centering
\caption{Temporal versus robust trust, $f=2$, 500 maps, with $95\%$ CIs on
recovery. ``nh'' is temporal no-harm. Recall shown in parentheses.}
\label{tab:temporal}
\begin{tabular}{llccccc}
\toprule
Attack & $\sigma$ & Base & Attack & Robust rec.\,(R) & Temporal rec.\,(R) & nh\\
\midrule
\multirow{4}{*}{Wall}
 & $0.0$ & $86.0$ & $71.8$ & $+13.9$ ($0.98$) & $+13.9$ ($0.98$) & $86.1$\\
 & $0.2$ & $79.1$ & $64.5$ & $+13.1$ ($0.90$) & $+14.1$ ($0.92$) & $78.8$\\
 & $0.4$ & $65.0$ & $54.8$ & $+7.8$ ($0.57$)  & $+9.2$ ($0.82$)  & $65.1$\\
 & $0.6$ & $53.4$ & $43.5$ & $+3.7$ ($0.26$)  & $+9.7$ ($0.70$)  & $53.1$\\
\midrule
\multirow{4}{*}{Camo.}
 & $0.0$ & $86.0$ & $63.5$ & $+20.3$ ($0.97$) & $+20.3$ ($0.97$) & $86.1$\\
 & $0.2$ & $79.1$ & $58.3$ & $+17.8$ ($0.86$) & $+18.6$ ($0.94$) & $78.8$\\
 & $0.4$ & $64.9$ & $47.7$ & $+11.1$ ($0.45$) & $+16.4$ ($0.81$) & $65.1$\\
 & $0.6$ & $53.4$ & $37.6$ & $+3.4$ ($0.13$)  & $\mathbf{+12.2}$ ($0.69$) & $53.0$\\
\bottomrule
\end{tabular}
\end{table}

The effect not only survives but strengthens as traitors are added, and it
does so well past the point where honest robots lose their local majority.
Table~\ref{tab:headline} isolates the worst case --- $\sigma=0.6$,
camouflage --- across traitor counts from one to seven of ten robots.
Because each honest robot is excluded from its own neighbour set, $f$
traitors leave it with at most $9-f$ honest neighbours: at $f=5$ an honest
robot's neighbourhood is already, in expectation, majority traitor, and by
$f=7$ only two of its nine potential neighbours are honest. The temporal
filter's recovery nonetheless climbs to a plateau of roughly $+15$ pp at
$f=4$--$6$ --- $+15.2$ pp (CI $[+11.9,+18.4]$) at $f=5$ and $+15.1$ pp
(CI $[+11.6,+18.6]$) at $f=6$ --- and remains $+10.9$ pp
(CI $[+7.4,+14.5]$) at $f=7$; every one of these intervals excludes zero.
The single-frame robust filter, by contrast, stays pinned near its floor and
is no longer statistically distinguishable from no defense at $f=7$
($+1.9$ pp, CI $[-1.3,+5.2]$). Temporal detection precision \emph{improves}
monotonically with the threat ($0.68\to0.82\to0.89\to0.93\to0.96\to0.97
\to0.98$): more traitors mean more genuine phantoms among the flagged links,
so the defense is most precise exactly when it matters most, while no-harm
stays flat at $\approx\!-0.4$ pp throughout.

This is the empirical content of operating without an honest local majority.
The verdict is pairwise --- each robot tests each neighbour against its own
sensing and never counts votes --- so the mechanism is structurally
independent of how many neighbours lie, and the data bear this out: the
defense's contribution is stable and significant whether honest robots are a
majority, a tie, or a minority. We are careful about what this does
\emph{not} claim. Absolute navigation success still falls as traitors are
added (undefended success drops from $44.4\%$ at $f=1$ to $\sim\!31\%$ at
$f=7$), because more liars do more damage; the temporal filter recovers a
stable, significant \emph{fraction} of that damage, it does not render
navigation traitor-count-invariant.

\begin{table}[t]\centering
\caption{Worst case ($\sigma=0.6$, camouflage) across traitor counts,
spanning the honest-majority boundary. Because each honest robot is excluded
from its own neighbour set, its neighbourhood is majority-traitor in
expectation once $f\ge5$. Recovery in pp; detection recall in parentheses;
temporal precision and no-harm in the last columns. 500 maps.}
\label{tab:headline}
\begin{tabular}{lccccc}
\toprule
$f$ & Undefended & Robust rec.\,(R) & Temporal rec.\,(R) & Prec. & No-harm\\
\midrule
$1$ & $44.4$ & $+1.9$ ($0.13$) & $+7.1$ ($0.69$)  & $0.68$ & $-0.4$\\
$2$ & $37.6$ & $+3.4$ ($0.13$) & $+12.2$ ($0.69$) & $0.82$ & $-0.4$\\
$3$ & $35.3$ & $+5.3$ ($0.12$) & $+13.6$ ($0.68$) & $0.89$ & $-0.3$\\
$4$ & $34.0$ & $+3.5$ ($0.12$) & $+14.7$ ($0.66$) & $0.93$ & $-0.3$\\
$5$ & $32.6$ & $+4.0$ ($0.13$) & $+15.2$ ($0.67$) & $0.96$ & $-0.4$\\
$6$ & $30.1$ & $+4.6$ ($0.13$) & $+15.1$ ($0.67$) & $0.97$ & $-0.4$\\
$7$ & $31.1$ & $+1.9$ ($0.13$) & $\mathbf{+10.9}$ ($0.66$) & $0.98$ & $-0.3$\\
\bottomrule
\end{tabular}
\end{table}

Figure~\ref{fig:offsethist} shows why the mechanism works: the distribution
of the aggregated offset magnitude cleanly separates honest neighbours
(concentrated near zero, as $\sqrt{2}\sigma$ noise averages out) from
fabrications (a persistent bias well above the threshold $\eta$). A probe
of this separation, before any trust gating, reports an area under the ROC
curve of $0.99$ under oracle association and $0.85$--$0.88$ under realistic
nearest-obstacle association (the latter rising with the accumulation window
to $0.88$ at the $K_{\min}=20$ operating point), both over the full 500-map
suite.

\subsection{An adaptive attacker faces a stealth/harm bind}\label{sec:res-adaptive}
A reasonable objection is that a static attacker understates the threat. We
therefore evaluate an attacker that knows the temporal filter and tunes its
fabrication to evade it, via four knobs. The clearest is the phantom
centre-offset: how far the fabricated obstacle is displaced from the real
one it hugs. Table~\ref{tab:bind} ($f=2$, $\sigma=0.6$) shows the resulting
bind. When the phantom sits on the real obstacle (offset $0$) it is
invisible to the filter (recall $0.03$) but also harmless (it blocks no new
space; harm $-0.8$ pp). As the offset grows, harm rises --- but so does
detection recall, in lockstep. The attacker cannot be both stealthy and
harmful: evading the temporal test requires collapsing the phantom onto real
geometry, which defeats the purpose of the attack.

\begin{table}[t]\centering
\caption{Stealth/harm bind: adaptive phantom centre-offset, $f=2$,
$\sigma=0.6$, 500 maps. Harm and detection recall rise together.}
\label{tab:bind}
\begin{tabular}{lcccc}
\toprule
Offset (m) & Undefended & Harm & Temporal rec. & Recall\\
\midrule
$0.0$ & $55.7$ & $-0.8$ & $+0.6$ & $0.03$\\
$0.5$ & $53.7$ & $+1.2$ & $+2.1$ & $0.07$\\
$1.0$ & $49.9$ & $+5.0$ & $+1.6$ & $0.48$\\
$1.5$ & $44.9$ & $+10.0$& $+5.6$ & $0.69$\\
$2.0$ & $42.1$ & $+12.8$& $+9.1$ & $0.70$\\
$2.5$ & $40.6$ & $+14.3$& $+8.5$ & $0.70$\\
\bottomrule
\end{tabular}
\end{table}

The bind holds across the noise range: sweeping the offset at
$\sigma\in\{0,0.2,0.4\}$ and $f\in\{1,2,3\}$ (nine configurations) yields
the same pattern in every one --- the harmless corner is undetected, the
harmful corner is caught, and detection recall at the harmful end stays high
($0.79$--$0.98$); the bind moreover persists through severe noise
($\sigma=0.6$, Table~\ref{tab:bind}), where harmful-end recall still holds at
$0.66$--$0.70$. The remaining three knobs are likewise self-defeating.
Per-frame jitter is zero-mean and so does not shift the temporal mean; it
only raises variance, and is self-defeating in practice: as the jitter grows,
the harm it inflicts falls sharply ($-6.8$ pp) while detection recall barely
moves ($0.70\to0.61$), leaving the defended success essentially unchanged.
Intermittent broadcasting (duty) dilutes the attacker's evidence but equally
dilutes its harm: at the lowest duty the harm falls to $+5.3$ pp while recall
falls to $0.41$. A surface gap large
enough to matter pushes the phantom centre outside the tolerance band, where
it is reliably caught. In no configuration does the adaptive attacker obtain
both meaningful harm and evasion.

\subsection{The navigation perception limit}\label{sec:res-limit}
Finally we distinguish the security result from an orthogonal limit. As
noise grows the attack-free base success itself falls, from $\sim\!86\%$ at
$\sigma=0$ to $53.4\%$ at $\sigma=0.6$. This is not a defense failure but a
navigation one: fine-tuning the base navigator under sensor-noise domain
randomization recovers only $\sim\!2.5$ pp at $\sigma=0.6$, indicating a
genuine perception-information limit --- position uncertainty comparable to
obstacle size cannot be undone by the policy. Crucially, the defense never
makes this worse (no-harm $\approx0$ at every level), and the two limits are
independent: the navigation limit bounds the ceiling, while the temporal
defense recovers the security gap up to that ceiling.
